% Dot batch 86; source SHA-256: 9198d175c5527bbc290626eba5b557948649678e5483f3f4c4f6b40129b3aad4
% Archived source: dot_updates_20261124/batches/batch086/EGA第五卷_格罗滕迪克致迪厄多内书信.tex
% Original body lines: 22--68; integration map: integration/late_integration.json
\begingroup
\def\newscope{V-L-letter}
\def\depth{\operatorname{depth}}
\def\FF{\mathcal F}
\head{第四篇}
\egasection{1. 致迪厄多内的信（1965年9月29日）}
\noindent 亲爱的迪厄多内：

谢谢你24日的来信，也谢谢你寄来的第16至第19节的目录。方便的时候，
我也很想收到第20、21节的暂定目录。把这两节收入第四章第四分册，我同意。
不过，你打算怎样拆分我原来的第20节？拆成两部分后，各叫什么题目呢？\cmt{L18-086-01}{“第四章第四分册”按法文 \emph{fascicule 4 du Chap IV} 译。原来的第20节是手写补入的宾语。英稿的“第四篇”“第1节”是编排标题，法文此页以日期和称呼开头。}

我已开始在整个计划里弄不清头绪了，而有时需要引到某一节而不至于说得太离谱，
这总归方便。下面是我理解的现行计划，你看是不是这样：\cmt{L18-086-02}{英译注82提醒原处用了不甚文雅的口语。此处按法文俚语的语气译，不把其字面粗俗义误作技术内容。下列番号为信中暂拟目录，不是正式出版目录。}

\begin{list}{}{\setlength{\leftmargin}{3em}\setlength{\itemsep}{1pt}\setlength{\topsep}{4pt}\setlength{\parsep}{0pt}}
 \item[20.] ？？？
 \item[21.] ？？？
 \item[22.] 线性系；Picard 群的补充。
 \item[23.] 格拉斯曼概形。
 \item[24.] 光滑型；普通二次奇点。
 \item[25.] 超平面截面，还有那一堆乱七八糟的东西。\cmt{L18-086-03}{保留 \emph{et bordel} 的戏谑口气。英译注83同样指出此处用语不雅，并给出一段残缺的英释；不将其另造为目录条目。}
 \item[26.] 结式与判别式。
 \item[27.] 无穷小扩张。
\end{list}

\begin{samepage}
第25节大概会很长，我已经能想象你打算把它再拆成两节了！
不过，\(27=3^3\) 可真是个漂亮的数。\cmt{L18-086-04}{法148与英112的页面上均有清楚的上标\(3\)；英稿纯文本抽取会把它排成“33”，不是英译本身的算式错误。}
\par\end{samepage}

第18节原先那篇附录，绝不能以我的名义发表。
你的写法，与我交给你的那些含混的手稿笔记几乎已没什么共同之处；
说不定我根本没给过你什么笔记，只对你说了句：“照完备环的办法做就行了……”
另一方面，若让可能用得上它的人看不到你这番整理的成果，那也可惜
（总会有人用得上的……）。所以，请你再考虑一下，把它写成一篇合著文章。\cmt{L18-086-05}{依照法148末至法149首的连续句读。英112在“即便我给过你”附近拆句并留有 \emph{que} 残字；原信的意思是作者甚至不确定自己实际交过笔记。此处只译合著建议，不据此判定后来发表时的实际署名。}

\begin{flushright}
祝好！\\[3pt]
A.~格罗滕迪克
\end{flushright}

\egasection{信末划删段落（另录）}
\noindent\emph{原在合著建议之后、落款之前。}\cmt{L18-086-06}{法149可见一道长斜线划过整个段落，落款不在划线上。英译注84也将划去的范围标为从合著建议后的段首到落款之前；英稿把末尾法语的意思再译了一次，不作为原信重复句。}

关于第20节的10.9.1，当然要用到这样一个事实：
给定 \(n\)，在 \(Z_\lambda\) 中，\(\FF_\lambda\) 限制到纤维后在该点的深度
\(>n\) 的点组成一个可构造集。（在我们所作的平坦和有限表示假设下，
第12节甚至应该已证明这个集合是开的。）
由于它在 \(Z\) 中的逆像是全体，到了稍晚于 \(\lambda\) 的某一阶段，那里也已经是全体了。
归根结底，反复出现的还真是同一个论证！\cmt{L18-086-07}{法149划线和比较号接近，放大可辨比较号为 \(>\)，英112也作 \(>\)。这个阶段结论按书信所指的极限逼近语境译出；本信没有重列该处的完整假设，也没有把括号内的开性说明表述为无条件断言。}
\par\endgroup
